\chapter{A Concrete Naming Certificate for $\MetaCICParallelDiamondFrontierUp$}
\label{ch:concrete-instances-metacic-parallel-diamond-frontier-namecert}
\origin{ai}

The $\MetaCICParallelDiamondFrontierUp$ packet realizes the residual-substitution frontier of \autoref{ch:metacic-confluence-audit-map} as finite confluence-facing NameCert data. It sits beside the residual-substitution boundary of \autoref{ch:concrete-instances-metacic-residual-substitution-boundary-namecert}, the confluence critical-pair boundary of \autoref{ch:concrete-instances-metacic-confluence-critical-pair-boundary-namecert}, the closed checker seal of \autoref{ch:concrete-instances-closed-type-checker-seal-namecert}, and the closed-normal confluence fragment of \autoref{ch:concrete-instances-metacicconfluencenormalfragment-namecert}. The packet records the conditional parallel-diamond frontier without claiming residual-substitution compatibility, full closed Church--Rosser, subject reduction, evaluator equality, or kernel consistency.

The retained $P,K,J,R,S,C,B,O$ rows are consumed substantively by the MetaCIC critical-path packet of \autoref{ch:concrete-instances-metaciccriticalpath-namecert}, the confluence audit packet of \autoref{ch:concrete-instances-metacic-confluence-audit-packet-namecert}, and the normalization-confluence bridge route of \autoref{ch:concrete-instances-metacic-normalization-confluence-bridge-namecert}. The critical-path reader consumes the residual-premise and obstruction rows, while the confluence audit packet consumes the candidate-join and bounded-check rows before any confluence consumer may use the frontier. The route exports no full Church--Rosser theorem and no subject-reduction theorem.

\falsifiablePrediction{Every accepted $\MetaCICParallelDiamondFrontierUp$ confluence-facing read exposes the residual-substitution premise row before any parallel-diamond handoff can be consumed.}

\independenceWitness{metacic-residual-substitution-boundary, metacic-confluence-critical-pair-boundary, closed-type-checker-seal, metacicconfluencenormalfragment: the residual-substitution boundary names the missing compatibility input without the parallel-diamond handoff; the critical-pair boundary records local peak data without the closed checker seal; the closed type-checker seal records bounded checking without confluence closure; and the closed-normal fragment records unconditional normal-source joins without the residual frontier. None is bijective to $\MetaCICParallelDiamondFrontierUp$.}

\begin{definition}[MetaCIC parallel-diamond frontier carrier]
\label{def:metacic-parallel-diamond-frontier-carrier}
A $\MetaCICParallelDiamondFrontierUp$ carrier is a finite $\BHist$ packet
$$
F=(P,K,J,R,S,C,B,O,H,T,G,N)
$$
with the following displayed rows:
\begin{enumerate}
\item $P$ is the residual-substitution premise row, naming the substitution compatibility input required by the parallel-diamond route.
\item $K$ is the critical-pair row, recording the local peak and its two parallel-reduction observations.
\item $J$ is the candidate-join row, recording the proposed common reduct and branch-local route requests.
\item $R$ is the residual-frontier row, listing residual descendants, binder-depth exposure, and branch-local substitution windows.
\item $S$ is the closed type-checker seal row consumed before a closed confluence-facing read can use the checker data.
\item $C$ is the closed-normal fragment row, recording the already available normal-source and atom-source confluence fragments.
\item $B$ is the bounded-check row over closedness, support, branch compatibility, and finished-normalization boundaries.
\item $O$ is the obstruction row naming the endpoints not supplied by the packet: residual-substitution discharge, full parallel diamond, subject reduction, arbitrary closed Church--Rosser, evaluator equality, and kernel consistency.
\item $H$ records componentwise $\hsame$ transport for premise, peak, join, residual, checker, fragment, bounded-check, obstruction, replay, provenance, and local naming rows.
\item $T$ records displayed $\Cont$ replay from a confluence query through $P,K,J,R,S,C,B,O,H,G,N$.
\item $G$ is the $\Pkg$ provenance over $\MetaCICParallelDiamondFrontierUp$, $\MetaCICResidualSubstitutionBoundaryUp$, $\MetaCICConfluenceCriticalPairBoundaryUp$, $\ClosedTypeCheckerSealUp$, $\MetaCICConfluenceNormalFragmentUp$, $\BHist$, $\hsame$, $\Cont$, and $\NameCert$.
\item $N$ is the local $\NameCert_{\MetaCICParallelDiamondFrontierUp}$ row.
\end{enumerate}
The classifier compares accepted packets only through $P,K,J,R,S,C,B,O,H,T,G,N$. It has no coordinate for a premise-free parallel diamond, subject reduction, arbitrary closed Church--Rosser, unrestricted beta-conversion decidability, evaluator equality, quotient equality of terms, or kernel consistency.
\end{definition}

\begin{theorem}[MetaCIC parallel-diamond frontier obligations]
\label{thm:metacic-parallel-diamond-frontier-obligations}
For every accepted $\MetaCICParallelDiamondFrontierUp$ carrier
$$
F=(P,K,J,R,S,C,B,O,H,T,G,N),
$$
the local $\NameCert_{\MetaCICParallelDiamondFrontierUp}$ surface consists exactly of residual-substitution premise retention through $P$, critical-pair exposure through $K$, candidate-join exposure through $J$, residual-frontier exposure through $R$, closed-checker handoff through $S$, closed-normal fragment exposure through $C$, bounded-check exposure through $B$, obstruction retention through $O$, componentwise transport through $H$, displayed replay through $T$, provenance through $G$, and local naming through $N$.
\end{theorem}
\leanchecked{BEDC.Derived.MetaCICParallelDiamondFrontierUp.MetacicParallelDiamondFrontierObligations}
\begin{proof}
Project the carrier of \autoref{def:metacic-parallel-diamond-frontier-carrier}. The confluence-facing rows are $P,K,J,R,S,C,B,O$, and the structural certificate rows are $H,T,G,N$.

The residual-substitution premise row remains first-class. Row $P$ is not discharged by the carrier; it is the named input that the conditional parallel-diamond route must still expose.

The local geometry is then displayed through $K,J,R$. The critical-pair, candidate-join, and residual-frontier rows provide finite peak data and branch-local residual windows, while the closed checker and closed-normal rows show which existing fragments may be consumed.

Finally $B,O,H,T,G,N$ keep the bounded checks, obstruction endpoints, transport, replay, provenance, and local naming attached to the same packet. Since no omitted theorem appears as a carrier coordinate, the listed rows exhaust the frontier obligations.
\end{proof}

\begin{theorem}[MetaCIC parallel-diamond frontier handoff]
\label{thm:metacic-parallel-diamond-frontier-handoff}
Every confluence consumer of an accepted $\MetaCICParallelDiamondFrontierUp$ packet must expose the residual-substitution premise row $P$, the critical-pair row $K$, the candidate-join row $J$, the residual-frontier row $R$, the closed checker seal row $S$, the closed-normal fragment row $C$, and the bounded-check row $B$ before any parallel-diamond handoff is read. The handoff may feed the conditional route recorded in \autoref{ch:metacic-confluence-audit-map}, but it does not discharge residual substitution, prove subject reduction, prove arbitrary closed confluence, decide evaluator equality, or establish kernel consistency.
\end{theorem}
\begin{proof}
Begin with the obligation surface of \autoref{thm:metacic-parallel-diamond-frontier-obligations}. It forces every confluence read through the named premise, peak, join, residual, checker, fragment, bounded-check, and obstruction rows.

The sibling packets preserve that boundary. The residual-substitution boundary supplies the named premise, the critical-pair boundary supplies local peak data, the closed type-checker seal supplies bounded checker data, and the normal fragment supplies only the closed-normal and atom-source cases already available.

The obstruction row $O$ remains part of the packet. Therefore the handoff is a finite frontier certificate for a conditional route, not a proof of the missing residual-substitution input or any stronger confluence endpoint.
\end{proof}

\begin{theorem}[MetaCIC parallel-diamond critical-pair exhaustion]
\label{thm:metacic-parallel-diamond-critical-pair-exhaustion}
For every accepted $\MetaCICParallelDiamondFrontierUp$ carrier
$$
F=(P,K,J,R,S,C,B,O,H,T,G,N),
$$
each residual-substitution critical pair that a confluence consumer may inspect is exposed by the displayed critical-pair row $K$ before the candidate-join row $J$ can be read. The read must retain the residual-substitution premise row $P$, route the branch-local residual windows through $R$, pass the closed checker seal $S$ and closed-normal fragment row $C$, and keep the bounded-check and obstruction rows $B,O$ visible. It exports no unlisted peak, hidden substitution lemma, candidate join detached from $K,R$, or closed-checker result outside the displayed frontier packet.
\end{theorem}
\begin{proof}
Project the carrier of \autoref{def:metacic-parallel-diamond-frontier-carrier}. The only rows that can mention residual substitution, critical pairs, candidate joins, residual windows, closed checking, and obstruction are $P,K,J,R,S,C,B,O$.

The residual-substitution premise is retained first. Row $P$ is not discharged by the packet, so every critical-pair read remains conditional on the same displayed premise.

The critical-pair exposure is then forced through $K$. A consumer cannot read $J$ as a candidate join until the peak and its two parallel-reduction observations have already been displayed in $K$ and routed through the branch-local residual windows in $R$.

Finally the closed-checker handoff is bounded by $S,C,B,O$. The checker seal, closed-normal fragment, bounded-check row, and obstruction row remain visible through the replay, provenance, and local naming rows $T,G,N$, so the packet exhausts the accepted critical pairs without adding any hidden confluence endpoint.
\end{proof}

\begin{theorem}[MetaCIC parallel-diamond frontier obstruction triad]
\label{thm:metacic-parallel-diamond-frontier-obstruction-triad}
Every accepted $\MetaCICParallelDiamondFrontierUp$ packet retains three independent obstruction classes throughout the conditional local-diamond read: residual-substitution premise discharge remains confined to the row $P$, subject-reduction compatibility remains outside the closed-checker and closed-normal rows $S,C$, and the full closed Church--Rosser, evaluator equality, quotient term equality, and kernel-consistency endpoints remain recorded only as refused endpoints in the obstruction row $O$. These obstructions are transported and replayed through $H,T,G,N$ without being collapsed into one another or promoted to completed confluence data.
\end{theorem}
\begin{proof}
Project the carrier rows $P,O,B,S,C$ from \autoref{def:metacic-parallel-diamond-frontier-carrier}. Row $P$ is the residual-substitution premise, row $O$ stores the refused endpoints, and rows $B,S,C$ expose only bounded checking, closed checking, and closed-normal fragment data.

Apply the frontier handoff theorem and critical-pair exhaustion, \autoref{thm:metacic-parallel-diamond-frontier-handoff} and \autoref{thm:metacic-parallel-diamond-critical-pair-exhaustion}. They keep the local peak under the displayed premise and force the closed-checker route to retain $S,C,B,O$.

Finally use the route theorem \autoref{thm:metacic-parallel-diamond-frontier-route}. Transport, replay, provenance, and local naming preserve the same rows rather than discharging $P$, deriving subject reduction from $S,C$, or turning the obstruction row into Church--Rosser, evaluator equality, quotient equality, or kernel consistency.
\end{proof}

\begin{theorem}[MetaCIC critical-pair obstruction budget]
\label{thm:metacic-parallel-diamond-frontier-critical-pair-obstruction-budget}
Every critical-pair read consumed by a $\MetaCICParallelDiamondFrontierUp$
parallel-diamond frontier must preserve the residual-substitution premise row
$P$ and the obstruction row $O$ through the displayed route
$$
P,K,J,R,S,C,B,O,H,T,G,N .
$$
In particular, a consumer may not retype the handoff as premise-free
confluence, subject reduction, evaluator equality, quotient term equality, or
kernel consistency.
\end{theorem}
\begin{proof}
Start with the obligation surface of
\autoref{thm:metacic-parallel-diamond-frontier-obligations}. It makes
$P,K,J,R,S,C,B,O,H,T,G,N$ the carrier rows of the accepted frontier packet,
with $P$ and $O$ displayed as ordinary rows rather than external commentary.

Apply the route theorem
\autoref{thm:metacic-parallel-diamond-frontier-route}. Any confluence-audit
consumer must expose the critical-pair row $K$, candidate row $J$,
residual-frontier row $R$, checker rows $S,C,B$, obstruction row $O$, and
structural rows $H,T,G,N$ before using the conditional handoff.

Use premise retention,
\autoref{thm:metacic-parallel-diamond-frontier-premise-retention}. The
residual-substitution premise row $P$ remains visible through the critical
pair, residual frontier, closed checker, and closed-normal rows whenever the
local-diamond handoff is consumed.

Finally apply discharge obstruction,
\autoref{thm:metacic-parallel-diamond-frontier-discharge-obstruction}. Since
no separate residual-substitution compatibility certificate is supplied by
this packet, the obstruction row $O$ keeps the refused confluence, typing,
evaluator, quotient, and kernel endpoints outside the critical-pair budget.
\end{proof}

\begin{theorem}[MetaCIC residual peak scope]
\label{thm:metacic-parallel-diamond-frontier-residual-peak-scope}
Every candidate-mediated residual peak read from an accepted $\MetaCICParallelDiamondFrontierUp$ packet must expose the MetaCIC critical-pair envelope, residual substitution dependency, normalization candidate scope, and closure-trace rows before a diamond frontier read is public. For a carrier
$$
F=(P,K,J,R,S,C,B,O,H,T,G,N),
$$
the visible scope is
$$
P,K,J,R,S,C,B,O,H,T,G,N,
$$
where $K$ records the critical-pair envelope, $P$ retains the residual-substitution dependency, $J,R,C$ carry the candidate-normalization and residual-frontier scope, and $H,T,G,N$ preserve the closure trace through transport, replay, provenance, and local naming. The read is a finite frontier boundary; it is not a confluence theorem, a strong-normalization theorem, subject reduction, evaluator equality, quotient term equality, or kernel consistency.
\end{theorem}
\begin{proof}
Project the carrier of \autoref{def:metacic-parallel-diamond-frontier-carrier}. The residual peak has no row outside $P,K,J,R,S,C,B,O,H,T,G,N$, so every public diamond-frontier read must be assembled from that finite spine.

Use critical-pair exhaustion. It displays the local peak through $K$ and prevents the candidate row $J$ from being read before the residual-substitution premise $P$ and branch-local residual row $R$ remain visible.

Next apply the frontier handoff and residual-substitution scope routes already present in the chapter. These keep the candidate-normalization surface inside $J,R,S,C,B,O$ and retain the obstruction row rather than closing the missing residual-substitution dependency.

Finally read the closure trace through $H,T,G,N$. Transport, replay, provenance, and local naming preserve the same candidate-mediated residual peak for downstream consumers, but none of those rows supplies confluence, strong normalization, subject reduction, evaluator equality, quotient equality, or kernel consistency.
\end{proof}

\begin{theorem}[MetaCIC parallel-diamond frontier route]
\label{thm:metacic-parallel-diamond-frontier-route}
Every confluence-audit consumer that reads an accepted $\MetaCICParallelDiamondFrontierUp$ packet must expose the residual-substitution premise row $P$, critical-pair row $K$, candidate-join row $J$, residual-frontier row $R$, closed-checker seal $S$, closed-normal fragment $C$, bounded-check row $B$, obstruction row $O$, transport row $H$, replay row $T$, provenance row $G$, and local name row $N$ before using a conditional parallel-diamond handoff. The route does not turn the residual-substitution premise into a theorem, does not prove full Church--Rosser, and does not discharge subject reduction, evaluator equality, quotient term equality, or kernel consistency.
\end{theorem}
\begin{proof}
Begin with the carrier of \autoref{def:metacic-parallel-diamond-frontier-carrier}. It displays exactly the premise, peak, join, residual, checker, fragment, bounded-check, obstruction, transport, replay, provenance, and local naming rows.

The obligation theorem \autoref{thm:metacic-parallel-diamond-frontier-obligations} keeps $P,K,J,R,S,C,B,O$ as first-class rows. A confluence-audit consumer therefore cannot use the handoff before the missing residual-substitution premise and the obstruction boundary are visible.

The handoff theorem \autoref{thm:metacic-parallel-diamond-frontier-handoff} supplies only a conditional frontier route. The added transport, replay, provenance, and local name rows preserve that boundary for downstream audit consumers.

Since no row discharges residual substitution or full closed confluence, the route is a sharper obligation boundary rather than a completed Church--Rosser certificate.
\end{proof}

\begin{theorem}[MetaCIC local-diamond frontier]
\label{thm:metacic-local-diamond-frontier}
Let
$$
F_0=(P,K,J_0,R,S,C,B,O,H,T,G,N)
\quad\text{and}\quad
F_1=(P,K,J_1,R,S,C,B,O,H,T,G,N)
$$
be accepted $\MetaCICParallelDiamondFrontierUp$ packets with the same residual-substitution premise row $P$ and the same critical-pair row $K$. If the candidate-join rows $J_0$ and $J_1$ agree through the displayed residual-frontier row $R$, then every confluence-audit consumer reads the same conditional local-diamond handoff through $S,C,B,O,H,T,G,N$. The obstruction row $O$ remains visible in the handoff, so the result does not discharge residual-substitution compatibility, subject reduction, full closed Church--Rosser, evaluator equality, quotient equality of terms, or kernel consistency.
\end{theorem}
\begin{proof}
Project the two carriers through \autoref{def:metacic-parallel-diamond-frontier-carrier}. The hypotheses identify the premise and peak rows $P,K$, and they make the two candidate-join rows agree at the residual-frontier row $R$.

Apply \autoref{thm:metacic-parallel-diamond-frontier-obligations} to each packet. The only confluence-facing rows available to a consumer before replay are $P,K,J,R,S,C,B,O$, with the structural rows $H,T,G,N$ attached to the same packet.

Now use \autoref{thm:metacic-parallel-diamond-frontier-handoff} and \autoref{thm:metacic-parallel-diamond-frontier-route}. Since $P,K$ agree and $J_0,J_1$ agree through $R$, both handoff readings reach the same closed-checker row $S$, closed-normal fragment row $C$, bounded-check row $B$, obstruction row $O$, transport row $H$, replay row $T$, provenance row $G$, and local name row $N$.

The obstruction row is not removed by this comparison. It is carried in both packets and remains part of the consumer-visible handoff, so the theorem records a local-diamond frontier fact rather than a proof of the missing residual-substitution or closed confluence endpoints.
\end{proof}

\begin{theorem}[MetaCIC parallel-diamond residual square]
\label{thm:metacic-parallel-diamond-frontier-residual-square}
Every accepted $\MetaCICParallelDiamondFrontierUp$ carrier
$$
F=(P,K,J,R,S,C,B,O,H,T,G,N)
$$
exposes a residual square only through the retained residual-substitution premise row $P$, the critical-pair row $K$, the candidate-join row $J$, the residual-frontier row $R$, the closed-checker seal $S$, the closed-normal fragment $C$, the bounded-check row $B$, the obstruction row $O$, and the transport, replay, provenance, and local naming rows $H,T,G,N$. Any consumer of the square must therefore factor through those rows before using a conditional local-diamond handoff. The square does not discharge residual substitution, subject reduction, full closed Church--Rosser, evaluator equality, quotient term equality, or kernel consistency.
\end{theorem}
\begin{proof}
Project the carrier through \autoref{def:metacic-parallel-diamond-frontier-carrier}. Its confluence-facing rows are exactly $P,K,J,R,S,C,B,O$, and its structural rows are exactly $H,T,G,N$.

Use critical-pair exhaustion, \autoref{thm:metacic-parallel-diamond-critical-pair-exhaustion}. It places the local peak data inside the displayed critical-pair row $K$ and keeps the closed-checker and obstruction rows visible.

Apply the local-diamond frontier theorem \autoref{thm:metacic-local-diamond-frontier}. A conditional local-diamond read can compare candidate joins only through the residual-frontier row $R$, with $S,C,B,O,H,T,G,N$ still present in the handoff.

Finally route the resulting read through \autoref{thm:metacic-parallel-diamond-frontier-route}. Since that route retains $P,K,J,R,S,C,B,O,H,T,G,N$ before any consumer sees the handoff, the residual square is a frontier obligation over the displayed rows rather than any of the refused residual-substitution, typing, confluence, equality, quotient, or kernel endpoints.
\end{proof}

\begin{theorem}[MetaCIC parallel-diamond residual-substitution scope]
\label{thm:metacic-parallel-diamond-frontier-residual-substitution-scope}
Every use of the $\MetaCICParallelDiamondFrontierUp$ local-diamond frontier must pass through the residual-substitution boundary, the confluence critical-pair boundary, and the closed-normal fragment before it claims a local join. In carrier terms, a consumer must expose the residual-substitution premise row $P$, the critical-pair row $K$, the residual-frontier row $R$, the closed checker seal row $S$, and the closed-normal fragment row $C$ before reading the candidate-join handoff. The resulting claim remains a local frontier claim; it does not close full Church--Rosser, subject reduction, evaluator equality, quotient equality of terms, or kernel consistency.
\end{theorem}
\begin{proof}
Apply the route theorem \autoref{thm:metacic-parallel-diamond-frontier-route}. It requires the residual-substitution premise row, the critical-pair row, the residual-frontier row, the closed-checker seal, and the closed-normal fragment to remain visible before the conditional handoff is consumed.

Then apply the local-diamond frontier theorem \autoref{thm:metacic-local-diamond-frontier}. The agreement of candidate-join rows is read only through the displayed residual-frontier row, and the obstruction row remains visible in the handoff.

The sibling boundaries named in the chapter introduction supply exactly the surrounding scope: the residual-substitution boundary names the missing compatibility premise, the confluence critical-pair boundary names the local peak surface, and the closed-normal fragment names the already available normal-source joins. Since none of these rows removes the obstruction row, the theorem records a residual-substitution-scoped local join rather than any stronger confluence endpoint.
\end{proof}

\begin{theorem}[MetaCIC parallel-diamond frontier premise retention]
\label{thm:metacic-parallel-diamond-frontier-premise-retention}
Any accepted $\MetaCICParallelDiamondFrontierUp$ packet whose local-diamond handoff is consumed retains the residual-substitution premise row $P$ through the critical-pair row $K$, residual-frontier row $R$, closed-checker row $S$, and closed-normal row $C$. No consumer may reinterpret that handoff as residual-substitution compatibility, premise-free parallel diamond, full Church--Rosser, subject reduction, evaluator equality, quotient equality of terms, or kernel consistency.
\end{theorem}
\begin{proof}
Apply the obligations theorem \autoref{thm:metacic-parallel-diamond-frontier-obligations}. It makes $P$ a first-class residual-substitution premise row and keeps $K,R,S,C$ among the confluence-facing rows rather than structural metadata.

Next use the handoff and route theorems, \autoref{thm:metacic-parallel-diamond-frontier-handoff} and \autoref{thm:metacic-parallel-diamond-frontier-route}. They require a consumer to expose the premise, critical-pair, residual-frontier, checker, and closed-normal rows before the conditional handoff is read.

Finally apply the residual-substitution scope theorem \autoref{thm:metacic-parallel-diamond-frontier-residual-substitution-scope}. Since the local-diamond frontier remains scoped by the residual-substitution boundary and the obstruction row is still visible, the retained handoff cannot be retyped as any of the refused confluence endpoints.
\end{proof}

\begin{theorem}[MetaCIC parallel-diamond closed-normal join socket]
\label{thm:metacic-parallel-diamond-frontier-closed-normal-join-socket}
Any closed-normal join socket consumed by a MetaCIC confluence audit must expose the residual-substitution premise row $P$, critical-pair row $K$, candidate-join row $J$, residual-frontier row $R$, closed-checker seal $S$, closed-normal fragment $C$, bounded-check row $B$, obstruction row $O$, componentwise $\hsame$ support row $H$, displayed $\Cont$ replay row $T$, $\Pkg$ provenance row $G$, and local $\NameCert$ row $N$ before any local join is read. The socket retains the residual-substitution premise and obstruction boundary; it does not claim residual-substitution compatibility or full closed Church--Rosser.
\end{theorem}
\begin{proof}
Start with the obligations and handoff theorems, \autoref{thm:metacic-parallel-diamond-frontier-obligations} and \autoref{thm:metacic-parallel-diamond-frontier-handoff}. They list $P,K,J,R,S,C,B,O,H,T,G,N$ as the confluence-facing packet and require the handoff to stay inside those rows.

Use the route and local-diamond theorems, \autoref{thm:metacic-parallel-diamond-frontier-route} and \autoref{thm:metacic-local-diamond-frontier}. These force any local join to pass through the critical-pair, candidate-join, residual-frontier, closed-checker, closed-normal, bounded-check, obstruction, and structural rows before a confluence audit reads it.

Finally apply residual-substitution scope and premise retention, \autoref{thm:metacic-parallel-diamond-frontier-residual-substitution-scope} and \autoref{thm:metacic-parallel-diamond-frontier-premise-retention}. The premise row $P$ remains visible throughout the closed-normal socket, and $O$ keeps the refused endpoints outside the handoff.
\end{proof}

\begin{theorem}[MetaCIC parallel-diamond frontier local-join ledger pullback]
\label{thm:metacic-parallel-diamond-frontier-local-join-ledger-pullback}
Let
$$
F=(P,K,J,R,S,C,B,O,H,T,G,N)
$$
be an accepted $\MetaCICParallelDiamondFrontierUp$ packet, and let
$$
L=(K,J,B,R,S,O,G_L,H_L,C_L,P_L,N_L)
$$
be an accepted $\MetaCICLocalJoinLedgerUp$ packet. Suppose the two packets share the same residual-substitution premise, critical-pair row, candidate-join row, bounded-check row, obstruction row, transport, replay, provenance, and local naming rows, with the frontier rows $R,S,C$ preserving the same residual-substitution scope as the ledger rows $R,S$. Then every confluence-audit consumer reads the local join through the same residual-substitution-scoped route. The pullback keeps the obstruction row visible, and it does not claim full Church--Rosser, subject reduction, evaluator equality, quotient term equality, or kernel consistency.
\end{theorem}
\begin{proof}
Project the frontier carrier through \autoref{def:metacic-parallel-diamond-frontier-carrier}. The accepted packet exposes $P,K,J,R,S,C,B,O,H,T,G,N$, with $P$ and $O$ retained as confluence-facing rows rather than external annotations.

Read the local-join ledger through \autoref{ch:concrete-instances-metacic-local-join-ledger-namecert}. The shared rows identify the local peak, candidate join, bounded checker, residual frontier, substitution boundary, obstruction readback, transport, replay, provenance, and local naming surfaces.

Apply the frontier obligations and route theorems, \autoref{thm:metacic-parallel-diamond-frontier-obligations} and \autoref{thm:metacic-parallel-diamond-frontier-route}. They force any confluence-audit consumer to expose the premise, critical-pair, candidate-join, residual-frontier, checker, bounded-check, obstruction, and structural rows before using a handoff.

Finally use premise retention, \autoref{thm:metacic-parallel-diamond-frontier-premise-retention}. Since the ledger rows are shared with the frontier rows and the obstruction row is preserved on both sides, the consumer-visible local join is read through the same residual-substitution-scoped route, not through an unconditional confluence endpoint.
\end{proof}

\begin{theorem}[MetaCIC parallel-diamond typed boundary]
\label{thm:metacic-parallel-diamond-typed-boundary}
Any accepted $\MetaCICParallelDiamondFrontierUp$ packet whose conditional local-diamond handoff is read under a typed-window premise must retain the typed window, residual-substitution premise row $P$, critical-pair row $K$, candidate-join row $J$, residual-frontier row $R$, closed-checker row $S$, closed-normal row $C$, bounded-check row $B$, and obstruction row $O$ before the handoff is consumed. The packet may not be cited as a premise-free parallel diamond, subject reduction, full Church--Rosser theorem, evaluator equality, quotient term equality, or kernel-consistency certificate.
\end{theorem}
\begin{proof}
Apply the obligations theorem \autoref{thm:metacic-parallel-diamond-frontier-obligations}. It lists the premise, critical-pair, candidate-join, residual-frontier, closed-checker, closed-normal, bounded-check, and obstruction rows as the confluence-facing surface of the accepted packet.

Use the route theorem \autoref{thm:metacic-parallel-diamond-frontier-route}. A typed-window consumer can read the conditional handoff only after those rows have been exposed together with transport, replay, provenance, and local naming.

Apply the residual-substitution scope theorem \autoref{thm:metacic-parallel-diamond-frontier-residual-substitution-scope}. The typed window is read through the same residual-substitution boundary, confluence critical-pair boundary, and closed-normal fragment that scope the local join.

Finally use premise retention, \autoref{thm:metacic-parallel-diamond-frontier-premise-retention}. The retained premise and obstruction rows prevent the handoff from being retyped as a premise-free diamond or any of the refused confluence, equality, or kernel endpoints.
\end{proof}

\begin{theorem}[MetaCIC parallel-diamond frontier typed-window handoff]
\label{thm:metacic-parallel-diamond-frontier-typed-window-handoff}
Any typed-window consumer of an accepted $\MetaCICParallelDiamondFrontierUp$ packet must read the full route
$$
P,K,J,R,S,C,B,O,H,T,G,N
$$
before feeding candidate-normalization or confluence-audit consumers. In this route $P$ remains the residual-substitution premise, $O$ remains the obstruction row, and the handoff cannot be consumed as a premise-free parallel diamond, full Church--Rosser theorem, subject reduction, evaluator equality, quotient equality of terms, or kernel-consistency certificate.
\end{theorem}
\begin{proof}
Project the carrier of \autoref{def:metacic-parallel-diamond-frontier-carrier}. The displayed packet contains the premise, critical-pair, candidate-join, residual-frontier, checker, fragment, bounded-check, obstruction, transport, replay, provenance, and local naming rows.

Apply the obligations theorem and the typed boundary theorem, \autoref{thm:metacic-parallel-diamond-frontier-obligations} and \autoref{thm:metacic-parallel-diamond-typed-boundary}. Together they force a typed-window consumer to expose $P,K,J,R,S,C,B,O$ before any conditional handoff is read.

Use premise retention and discharge obstruction, \autoref{thm:metacic-parallel-diamond-frontier-premise-retention} and \autoref{thm:metacic-parallel-diamond-frontier-discharge-obstruction}. The residual-substitution premise row $P$ is still a premise, and the obstruction row $O$ still blocks discharge unless a separate compatibility certificate is supplied.

Finally route the exported handoff through the structural rows $H,T,G,N$. Candidate-normalization and confluence-audit consumers therefore receive exactly the retained typed-window spine, not a completed parallel-diamond, Church--Rosser, subject-reduction, evaluator-equality, quotient-equality, or kernel-consistency endpoint.
\end{proof}

\begin{theorem}[MetaCIC parallel-diamond frontier discharge obstruction]
\label{thm:metacic-parallel-diamond-frontier-discharge-obstruction}
Any accepted $\MetaCICParallelDiamondFrontierUp$ packet consumed as a local-diamond handoff still exposes the residual-substitution premise row $P$ and obstruction row $O$. Unless a separate $\MetaCICResidualSubstitutionCompatibilityUp$ certificate is supplied through \autoref{ch:concrete-instances-metacic-residual-substitution-compatibility-namecert}, no consumer may read that handoff as a premise-free parallel diamond, full closed confluence, subject reduction, evaluator equality, quotient term equality, or kernel consistency.
\end{theorem}
\begin{proof}
Project the carrier of \autoref{def:metacic-parallel-diamond-frontier-carrier}. The rows $P$ and $O$ are displayed coordinates of the accepted packet, not annotations around the handoff.

Apply premise retention, \autoref{thm:metacic-parallel-diamond-frontier-premise-retention}. It keeps $P$ visible through the critical-pair, residual-frontier, closed-checker, and closed-normal rows whenever the local-diamond handoff is consumed.

Now use the sibling route theorem \autoref{thm:metacic-parallel-diamond-frontier-sibling-route}. The exported contribution is exactly retained premise and obstruction data for confluence-facing siblings, so the handoff cannot be consumed as a premise-free diamond.

The only admissible discharge surface is the residual-substitution compatibility chapter named in the statement. Without that separate certificate, the obstruction row $O$ continues to block full closed confluence, subject reduction, evaluator equality, quotient term equality, and kernel consistency.
\end{proof}

\begin{theorem}[MetaCIC parallel-diamond frontier compatibility discharge surface]
\label{thm:metacic-parallel-diamond-frontier-compatibility-discharge-surface}
For an accepted $\MetaCICParallelDiamondFrontierUp$ packet, only a separate $\MetaCICResidualSubstitutionCompatibilityUp$ row may turn the retained $P,O$ premise-obstruction surface into a discharged local diamond. Without that compatibility row, every confluence consumer retains the full conditional spine
$$
P,K,J,R,S,C,B,O,H,T,G,N
$$
and may read the packet only as a conditional parallel-diamond frontier, not as premise-free confluence, subject reduction, evaluator equality, quotient term equality, or kernel consistency.
\end{theorem}
\begin{proof}
Start with the discharge obstruction theorem, \autoref{thm:metacic-parallel-diamond-frontier-discharge-obstruction}. It identifies $P$ as the retained residual-substitution premise and $O$ as the obstruction row that remains visible unless the separate compatibility chapter supplies a certificate.

Next use premise retention, \autoref{thm:metacic-parallel-diamond-frontier-premise-retention}. The premise row is not confined to the statement of the theorem; it is transported through the critical-pair, residual-frontier, closed-checker, and closed-normal rows of the accepted carrier.

Finally apply the confluence handoff, \autoref{thm:metacic-parallel-diamond-confluence-handoff}. That handoff exports the complete route $P,K,J,R,S,C,B,O,H,T,G,N$ to confluence-facing consumers. Since no row in this packet is the separate $\MetaCICResidualSubstitutionCompatibilityUp$ certificate, the handoff remains conditional and the refused endpoints stay outside the frontier.
\end{proof}

\begin{theorem}[MetaCIC parallel-diamond closed-checker retention]
\label{thm:metacic-parallel-diamond-closed-checker-retention}
Any accepted $\MetaCICParallelDiamondFrontierUp$ packet whose local-diamond handoff retains the residual-substitution premise row $P$ must also retain the closed-checker seal row $S$ and the closed-normal fragment row $C$ before that handoff can be read. Thus the conditional handoff boundary is the joint row surface $P,S,C$ together with the displayed critical-pair, residual-frontier, bounded-check, obstruction, transport, replay, provenance, and local naming rows. It is not a proof of residual substitution, subject reduction, evaluator equality, quotient term equality, or kernel consistency.
\end{theorem}
\begin{proof}
Apply the obligations theorem \autoref{thm:metacic-parallel-diamond-frontier-obligations}. It lists $P,S,C$ as confluence-facing rows of the same accepted carrier, so the checker seal and closed-normal fragment are not optional metadata around the premise row.

Next apply the handoff and route theorems, \autoref{thm:metacic-parallel-diamond-frontier-handoff} and \autoref{thm:metacic-parallel-diamond-frontier-route}. Both require the closed-checker seal and closed-normal fragment to be exposed before any conditional parallel-diamond handoff is consumed.

Finally use premise retention, \autoref{thm:metacic-parallel-diamond-frontier-premise-retention}. Since the retained premise row is read through $K,R,S,C$ and the obstruction row remains visible, the closed-checker seal and closed-normal fragment remain part of the boundary rather than consequences of the handoff. The refused endpoints therefore stay outside the packet.
\end{proof}

\begin{theorem}[MetaCIC parallel-diamond frontier sibling route]
\label{thm:metacic-parallel-diamond-frontier-sibling-route}
Every accepted $\MetaCICParallelDiamondFrontierUp$ packet that exports its premise-retention handoff feeds two confluence-facing sibling routes: the candidate-normalization confluence handoff of \autoref{ch:concrete-instances-metacic-candidate-normalization-confluence-handoff-namecert} and the audit-packet boundary of \autoref{ch:concrete-instances-metacic-confluence-audit-packet-namecert}. The exported contribution is exactly the residual-substitution premise row $P$, the local critical-pair row $K$, the residual-frontier row $R$, and the obstruction row $O$ retained by \autoref{thm:metacic-parallel-diamond-frontier-premise-retention}; it does not supply a premise-free diamond, a completed residual-substitution theorem, full closed confluence, subject reduction, evaluator equality, quotient equality of terms, or kernel consistency.
\end{theorem}
\begin{proof}
Apply \autoref{thm:metacic-parallel-diamond-frontier-premise-retention}. Any consumed local-diamond handoff keeps the residual-substitution premise row visible through the critical-pair, residual-frontier, closed-checker, and closed-normal rows.

The candidate-normalization confluence handoff can therefore read this packet only as retained frontier data: $P$ names the missing residual-substitution input, $K$ and $R$ name the local peak and residual window, and $O$ keeps the obstruction endpoints visible before candidate-normalization routes are consumed.

The confluence audit packet reads the same rows as an audit boundary. It can cite the retained premise and obstruction rows when it records a conditional local-diamond route, but it cannot erase them into a completed confluence theorem.

Thus the sibling lattice is substantive rather than bibliographic: the frontier supplies exactly the residual-substitution premise retention and local-diamond obstruction rows consumed by those confluence-facing packets, while the refused endpoints remain outside the carrier.
\end{proof}

\begin{theorem}[MetaCIC parallel-diamond confluence handoff]
\label{thm:metacic-parallel-diamond-confluence-handoff}
An accepted $\MetaCICParallelDiamondFrontierUp$ packet whose local-diamond, premise-retention, and residual-substitution rows have been displayed exports exactly the conditional confluence handoff consumed by \autoref{ch:concrete-instances-metacic-candidate-normalization-confluence-handoff-namecert} and bounded by \autoref{ch:concrete-instances-parallel-diamond-residual-namecert}. In carrier terms, the handoff is the retained route through $P,K,J,R,S,C,B,O,H,T,G,N$: $P$ remains the residual-substitution premise, $K,J,R$ expose the local-diamond and residual windows, $S,C,B$ keep the closed checker and bounded fragment visible, and $O$ records the still-open endpoints. The handoff does not claim global confluence, full strong normalization, subject reduction, evaluator equality, quotient equality of terms, or kernel consistency.
\end{theorem}
\begin{proof}
Project the frontier carrier. The rows $P,K,J,R,S,C,B,O$ are the confluence-facing data, while $H,T,G,N$ transport, replay, source, and name the same packet.

Apply \autoref{thm:metacic-parallel-diamond-frontier-premise-retention}. The residual-substitution premise row $P$ is retained through the local critical-pair and residual-frontier rows, so the exported route cannot be read as a premise-free diamond.

Next use \autoref{thm:metacic-parallel-diamond-frontier-discharge-obstruction}. It keeps the obstruction row $O$ visible unless a separate residual-substitution compatibility certificate is supplied, and therefore preserves the boundary checked by the parallel-diamond residual packet.

Finally apply the sibling route theorem. The candidate-normalization handoff receives only the closed checker boundary, local-diamond window, residual ledger, retained premise, and obstruction data displayed by this packet. Since no carrier row contains global confluence, full strong normalization, subject reduction, evaluator equality, quotient equality, or kernel consistency, none of those endpoints is exported by the handoff.
\end{proof}

\begin{theorem}[MetaCIC parallel-diamond premise preservation]
\label{thm:metacic-parallel-diamond-frontier-premise-preservation}
Any confluence consumer that reads the conditional parallel-diamond handoff of \autoref{thm:metacic-parallel-diamond-confluence-handoff} must retain the residual-substitution premise row $P$ through the full packet spine
$$
P,K,J,R,S,C,B,O,H,T,G,N .
$$
In particular, the handoff cannot erase $P$ into a completed residual-substitution theorem, premise-free parallel diamond, full Church--Rosser theorem, subject-reduction theorem, evaluator equality, quotient equality of terms, or kernel-consistency certificate.
\end{theorem}
\begin{proof}
Start with \autoref{thm:metacic-parallel-diamond-frontier-obligations}. It makes the residual-substitution premise row $P$ part of the local $\NameCert_{\MetaCICParallelDiamondFrontierUp}$ surface together with $K,J,R,S,C,B,O,H,T,G,N$.

Next apply premise retention, \autoref{thm:metacic-parallel-diamond-frontier-premise-retention}. The retained premise is still visible when the local-diamond handoff is consumed, and the discharge-obstruction theorem \autoref{thm:metacic-parallel-diamond-frontier-discharge-obstruction} keeps $O$ displayed unless a separate compatibility certificate is supplied.

Finally use the confluence handoff theorem itself. Its exported route is exactly the conditional path through $P,K,J,R,S,C,B,O,H,T,G,N$, so any consumer that drops $P$ or reads the route as one of the refused endpoints has left the packet described by the carrier.
\end{proof}

\begin{theorem}[MetaCIC parallel-diamond residual budget]\label{thm:metacic-parallel-diamond-frontier-budget}
Any exported $\MetaCICParallelDiamondFrontierUp$ parallel-diamond budget that feeds candidate-normalization or audit-packet consumers must consume the retained residual-substitution premise row $P$, critical-pair row $K$, residual-frontier row $R$, closed-checker row $S$, closed-normal row $C$, obstruction row $O$, and structural rows $H,T,G,N$ before the consumer receives the handoff. The budget cannot omit $P$ or $O$, and it cannot be read as a premise-free parallel diamond, residual-substitution theorem, full Church--Rosser theorem, subject-reduction theorem, evaluator equality, quotient equality of terms, or kernel-consistency certificate.
\end{theorem}
\begin{proof}
Start from the carrier obligations \autoref{thm:metacic-parallel-diamond-frontier-obligations}. They keep the residual-substitution premise, critical-pair, residual-frontier, closed-checker, closed-normal, obstruction, transport, replay, provenance, and naming rows inside the accepted packet.

Use the route and premise-retention theorems, \autoref{thm:metacic-parallel-diamond-frontier-route} and \autoref{thm:metacic-parallel-diamond-frontier-premise-retention}. They require the premise row $P$ to remain visible when a local-diamond handoff is consumed.

Next apply discharge obstruction and the sibling route, \autoref{thm:metacic-parallel-diamond-frontier-discharge-obstruction} and \autoref{thm:metacic-parallel-diamond-frontier-sibling-route}. The candidate-normalization and audit-packet consumers receive only the retained premise, critical-pair, residual-frontier, obstruction, and structural rows named by the frontier.

Finally use premise preservation \autoref{thm:metacic-parallel-diamond-frontier-premise-preservation}. Since its full spine includes $P,K,J,R,S,C,B,O,H,T,G,N$, any exported budget that forgets $P$ or $O$ has left the displayed packet and cannot be consumed as a BEDC parallel-diamond frontier.
\end{proof}

\begin{theorem}[MetaCIC parallel-diamond frontier obstruction monotonicity]
\label{thm:metacic-parallel-diamond-frontier-obstruction-monotonicity}
Let $F=(P,K,J,R,S,C,B,O,H,T,G,N)$ be an accepted $\MetaCICParallelDiamondFrontierUp$ packet. Any refinement of $F$ that preserves the displayed rows $P,K,J,R,S,C,B,O,H,T,G,N$ may strengthen the finite evidence attached to those rows, but it cannot weaken the obstruction row $O$ or discharge the residual-substitution premise row $P$ unless it also supplies a separate residual-substitution compatibility certificate. Consequently the refined packet remains a conditional frontier and cannot be consumed as full Church--Rosser, subject reduction, evaluator equality, quotient equality of terms, or kernel consistency.
\end{theorem}
\begin{proof}
Start with the obligation surface of \autoref{thm:metacic-parallel-diamond-frontier-obligations}. The premise row $P$ and obstruction row $O$ are carrier coordinates of the accepted packet, not side assumptions that a refinement may silently erase.

Apply premise retention and discharge obstruction, \autoref{thm:metacic-parallel-diamond-frontier-premise-retention} and \autoref{thm:metacic-parallel-diamond-frontier-discharge-obstruction}. They force the local-diamond handoff to keep $P$ visible and keep $O$ active unless a separate compatibility certificate supplies the missing residual-substitution discharge.

Finally use the residual budget theorem \autoref{thm:metacic-parallel-diamond-frontier-budget}. A refinement that preserves $P,K,J,R,S,C,B,O,H,T,G,N$ preserves the same budget boundary, so it cannot convert retained obstruction data into a premise-free endpoint. The stronger claims listed in the statement therefore remain outside the packet.
\end{proof}

\begin{theorem}[MetaCIC parallel-diamond residual confluence]
\label{thm:metacic-parallel-diamond-frontier-residual-confluence}
Every confluence consumer of an accepted $\MetaCICParallelDiamondFrontierUp$ packet must factor its exported result through the full residual spine
$$
P,K,J,R,S,C,B,O,H,T,G,N .
$$
The exported result is conditional residual confluence with the residual-substitution premise retained; it is not a premise-free parallel diamond, subject reduction theorem, evaluator equality, quotient term equality, closed Church--Rosser theorem, or kernel-consistency certificate.
\end{theorem}
\begin{proof}
Start with the obligation surface of \autoref{thm:metacic-parallel-diamond-frontier-obligations}. It makes $P,K,J,R,S,C,B,O,H,T,G,N$ the visible rows of any accepted frontier packet, with $P$ retained as a residual-substitution premise.

Use the route theorem and premise retention, \autoref{thm:metacic-parallel-diamond-frontier-route} and \autoref{thm:metacic-parallel-diamond-frontier-premise-retention}. Together they force confluence consumers to expose the premise, peak, join, residual, checker, fragment, bounded-check, obstruction, transport, replay, provenance, and local name rows before the handoff is consumed.

Apply discharge obstruction and the confluence handoff theorem, \autoref{thm:metacic-parallel-diamond-frontier-discharge-obstruction} and \autoref{thm:metacic-parallel-diamond-confluence-handoff}. These results keep the obstruction row visible and identify the exported route as the retained conditional handoff consumed by confluence-facing siblings.

Finally use the residual budget theorem \autoref{thm:metacic-parallel-diamond-frontier-budget}. Since the budget cannot omit $P$ or $O$, any exported confluence result remains conditional residual confluence over the displayed spine rather than one of the refused premise-free, typing, equality, quotient, Church--Rosser, or kernel endpoints.
\end{proof}

\begin{theorem}[MetaCIC residual-join readback]
\label{thm:metacic-parallel-diamond-frontier-residual-join-readback}
For an accepted $\MetaCICParallelDiamondFrontierUp$ carrier $F=(P,K,J,R,S,C,B,O,H,T,G,N)$, every residual-join readback consumed by bounded conversion must display the closed substitution boundary rows, candidate-mediated strong-normalization rows, residual budget rows, typed-example boundary rows, and the four L10 source faces
$$
\DyadicRatCoreUp,\quad \StreamNameUp,\quad \RegSeqRatUp,\quad \RealUp
$$
before reading the handoff. In carrier terms, the readback must expose $P,K,J,R,S,C,B,O,H,T,G,N$ together with the typed-window and L10 source rows, and it cannot be consumed as full confluence, global beta decidability, closed subject reduction, quotient term equality, selected limits, or $\RealUp$ completeness.
\end{theorem}
\begin{proof}
Start with residual confluence, \autoref{thm:metacic-parallel-diamond-frontier-residual-confluence}. It forces a confluence consumer to keep the residual-substitution premise, local critical pair, candidate join, residual frontier, checker seal, closed-normal fragment, bounded-check row, obstruction row, and structural rows visible through the full spine $P,K,J,R,S,C,B,O,H,T,G,N$.

Next use the typed-window handoff of \autoref{thm:metacic-parallel-diamond-frontier-typed-window-handoff}. A bounded-conversion consumer may read the frontier only after the typed window has retained the same residual spine for candidate-normalization and confluence-audit consumers.

Premise preservation, \autoref{thm:metacic-parallel-diamond-frontier-premise-preservation}, keeps the residual-substitution row $P$ from disappearing into a completed diamond. The local-diamond and critical-pair routes of \autoref{thm:metacic-local-diamond-frontier} and \autoref{thm:metacic-parallel-diamond-frontier-route} supply the candidate join and residual budget rows without discharging the obstruction row $O$.

The typed-example boundary and the four L10 source faces are then read as source rows, not as semantic quotients. They expose $\DyadicRatCoreUp$, $\StreamNameUp$, $\RegSeqRatUp$, and $\RealUp$ before the bounded-conversion handoff is consumed, so selected limits and ambient real completeness are unavailable.

Combining these routes gives a finite residual-join readback frontier. Since every step retains the premise and obstruction rows, the result is not full confluence, global beta decidability, closed subject reduction, quotient term equality, selected limits, or $\RealUp$ completeness.
\end{proof}

\begin{theorem}[MetaCIC critical-overlap residual budget]
\label{thm:metacic-parallel-diamond-frontier-critical-overlap-budget}
For an accepted $\MetaCICParallelDiamondFrontierUp$ carrier $F=(P,K,J,R,S,C,B,O,H,T,G,N)$, every critical overlap between two local reductions is budgeted by the finite row surface
$$
P,K,J,R,S,C,B,O,H,T,G,N .
$$
The overlap budget must expose the residual-substitution premise row $P$, the critical-pair row $K$, the candidate-join row $J$, the residual-frontier row $R$, the closed checker seal $S$, the closed-normal fragment $C$, the bounded-check row $B$, and the obstruction row $O$ before any local join is read. It does not supply source-equivalence transport, abstract carrier equivalence, premise-free parallel diamond, subject reduction, evaluator equality, quotient term equality, or kernel consistency.
\end{theorem}
\begin{proof}
Start with the obligation surface of \autoref{thm:metacic-parallel-diamond-frontier-obligations}. It makes the premise, critical-pair, candidate-join, residual-frontier, checker, fragment, bounded-check, obstruction, transport, replay, provenance, and local naming rows the finite carrier of an accepted frontier packet.

Use the residual-substitution scope theorem \autoref{thm:metacic-parallel-diamond-frontier-residual-substitution-scope}. A local join may be read only after the residual-substitution boundary, critical-pair boundary, and closed-normal fragment have remained visible.

Now apply premise preservation and the residual budget theorem, \autoref{thm:metacic-parallel-diamond-frontier-premise-preservation} and \autoref{thm:metacic-parallel-diamond-frontier-budget}. These keep $P$ and $O$ inside the exported budget and prevent a critical overlap from being consumed as a premise-free diamond.

Finally use residual confluence, \autoref{thm:metacic-parallel-diamond-frontier-residual-confluence}. Since the confluence consumer factors through the same displayed spine, the critical-overlap budget is finite and row-local rather than a transport across source equivalence or abstract carrier equivalence.
\end{proof}

\begin{theorem}[MetaCIC candidate-mediated residual route]
\label{thm:metacic-parallel-diamond-frontier-candidate-mediated-residual-route}
Let two local reductions be recorded by the critical-pair row $K$ of an accepted $\MetaCICParallelDiamondFrontierUp$ packet $F=(P,K,J,R,S,C,B,O,H,T,G,N)$. If the candidate-join row $J$ is accepted through the displayed residual-frontier row $R$, then both local reductions route to the shared residual packet
$$
J,R,S,C,B,O,H,T,G,N
$$
under the retained premise $P$. The route is candidate-mediated: it identifies the common residual packet available to the frontier consumer, but it does not discharge residual-substitution compatibility, source-equivalence transport, subject reduction, full closed Church--Rosser, evaluator equality, quotient term equality, or kernel consistency.
\end{theorem}
\begin{proof}
Project the carrier in \autoref{def:metacic-parallel-diamond-frontier-carrier}. The critical-pair row $K$ records the two local reduction observations, and the candidate-join row $J$ records the proposed common reduct together with branch-local route requests.

Apply the obligations theorem \autoref{thm:metacic-parallel-diamond-frontier-obligations}. It keeps $J$ and $R$ as displayed rows of the same packet, so the candidate route is not an external equivalence between carriers.

Use the residual-substitution scope theorem and premise preservation, \autoref{thm:metacic-parallel-diamond-frontier-residual-substitution-scope} and \autoref{thm:metacic-parallel-diamond-frontier-premise-preservation}. The two local reductions may reach the candidate only while $P$ remains visible and while the residual-frontier row controls the branch windows.

Finally apply residual confluence and residual-join readback, \autoref{thm:metacic-parallel-diamond-frontier-residual-confluence} and \autoref{thm:metacic-parallel-diamond-frontier-residual-join-readback}. They identify the consumer-facing route as the shared residual packet through $J,R,S,C,B,O,H,T,G,N$, with the premise and obstruction still retained.
\end{proof}

\begin{definition}[MetaCIC parallel-diamond frontier finite budget]
\label{def:metacic-parallel-diamond-frontier-budget}
The finite budget for a $\MetaCICParallelDiamondFrontierUp$ read is the row surface
$$
\mathcal{B}_{\mathrm{frontier}}=(W,K,P,J,R,S,C,B,O,H,T,G,N).
$$
Here $W$ is the typed-window row, $K$ is the critical-pair row, $P$ is the residual-substitution boundary, $J$ is the local-join row, $R$ is the residual-frontier row, $S$ is the closed-checker seal, $C$ is the closed-normal fragment, $B$ is the bounded-checker row, $O$ is the obstruction row, $H$ is componentwise $\hsame$ transport, $T$ is displayed $\Cont$ replay, $G$ is $\Pkg$ provenance, and $N$ is the local $\NameCert_{\MetaCICParallelDiamondFrontierUp}$ row. The budget has no coordinate for premise-free confluence, subject reduction, evaluator equality, quotient term equality, global normalization, or kernel consistency.
\end{definition}

\begin{theorem}[MetaCIC typed-window local-join scope]
\label{thm:metacic-parallel-diamond-frontier-local-join-scope}
Under a typed MetaCIC window, an accepted $\MetaCICParallelDiamondFrontierUp$ packet permits only local joins whose two branches pass through the retained residual-substitution premise $P$, critical-pair row $K$, candidate-join row $J$, residual-frontier row $R$, closed checker seal $S$, closed-normal fragment $C$, and bounded-check row $B$. The local join scope is therefore the typed-window row surface
$$
P,K,J,R,S,C,B,O,H,T,G,N ,
$$
with $O$ still visible. It is not a source-equivalence transport theorem, an abstract carrier-equivalence theorem, a premise-free parallel diamond, subject reduction, global beta decidability, evaluator equality, quotient term equality, or kernel consistency.
\end{theorem}
\begin{proof}
Start with the typed boundary and typed-window handoff, \autoref{thm:metacic-parallel-diamond-typed-boundary} and \autoref{thm:metacic-parallel-diamond-frontier-typed-window-handoff}. They force a typed-window consumer to expose the same residual spine before candidate-normalization or confluence-audit consumers can read a handoff.

Use the critical-overlap budget of \autoref{thm:metacic-parallel-diamond-frontier-critical-overlap-budget}. It keeps every local overlap inside $P,K,J,R,S,C,B,O,H,T,G,N$ and blocks readings that omit the premise or obstruction rows.

Next apply the candidate-mediated route \autoref{thm:metacic-parallel-diamond-frontier-candidate-mediated-residual-route}. The two branches can meet only through the displayed candidate row and residual-frontier row, not through an ambient source-equivalence or carrier-equivalence transport.

Finally use residual-join readback, \autoref{thm:metacic-parallel-diamond-frontier-residual-join-readback}. The readback keeps bounded conversion, typed windows, and the L10 source faces visible before the handoff, so the permitted local join remains a scoped MetaCIC window claim rather than any stronger confluence, typing, equality, quotient, or kernel endpoint.
\end{proof}

\begin{theorem}[MetaCIC parallel-diamond frontier local join budget]
\label{thm:metacic-parallel-diamond-frontier-local-join}
Any local join read accepted by a $\MetaCICParallelDiamondFrontierUp$ packet consumes the finite budget of \autoref{def:metacic-parallel-diamond-frontier-budget} before it can reach a normalization-facing consumer. In particular, the typed-window row, critical-pair row, residual-substitution boundary, local-join row, residual-frontier row, closed-checker seal, bounded-checker row, obstruction row, transport, replay, provenance, and local naming remain visible at the handoff.
\end{theorem}
\begin{proof}
Start with the typed-window local-join scope of \autoref{thm:metacic-parallel-diamond-frontier-local-join-scope}. It forces any accepted local join to expose the residual-substitution premise, critical pair, candidate join, residual-frontier row, checker rows, and obstruction row.

Apply the finite budget definition \autoref{def:metacic-parallel-diamond-frontier-budget}. The typed-window row and structural rows $H,T,G,N$ are part of the same budget surface, so the local join cannot be consumed as a free-standing confluence fact.

Use the residual budget theorem \autoref{thm:metacic-parallel-diamond-frontier-budget}. It keeps the obstruction and residual-substitution rows visible before normalization consumers read the packet. Therefore the handoff remains a finite typed-window claim and does not become premise-free normalization or global confluence.
\end{proof}

\begin{theorem}[MetaCIC parallel-diamond frontier residual boundary]
\label{thm:metacic-parallel-diamond-frontier-residual-boundary}
Residual-substitution compatibility remains a visible boundary in every accepted $\MetaCICParallelDiamondFrontierUp$ local-join read. The row $P$ of \autoref{def:metacic-parallel-diamond-frontier-budget} is consumed before $J,R,S,C,B$ can be used, and it is not erased into an unconditional confluence theorem, subject-reduction theorem, evaluator equality, quotient term equality, global normalization statement, or kernel-consistency certificate.
\end{theorem}
\begin{proof}
Project the budget surface of \autoref{def:metacic-parallel-diamond-frontier-budget}. The residual-substitution boundary $P$ is an explicit coordinate, not a proof obligation hidden behind the local-join row.

Apply premise preservation and the residual-substitution scope theorem, \autoref{thm:metacic-parallel-diamond-frontier-premise-preservation} and \autoref{thm:metacic-parallel-diamond-frontier-residual-substitution-scope}. These keep $P$ visible whenever the critical-pair, residual-frontier, checker, and bounded-check rows are read.

Finally use the local join budget theorem \autoref{thm:metacic-parallel-diamond-frontier-local-join}. Since every normalization-facing read consumes the full finite budget first, no consumer can erase the residual boundary and present the packet as an unconditional confluence or kernel result.
\end{proof}

\begin{theorem}[MetaCIC parallel-diamond frontier obligation basis]
\label{thm:metacic-parallel-diamond-frontier-obligation-basis}
Every accepted $\MetaCICParallelDiamondFrontierUp$ packet presents a finite obligation basis for the typed-window frontier. The basis is the row surface
$$
P,K,J,R,S,C,B,O,H,T,G,N,
$$
where $P$ is the residual-substitution boundary, $K$ is the critical-pair envelope, $J$ is the local-join row, $R$ is the residual-frontier row, $S$ is the closed checker seal, $C$ is the closed-normal fragment, $B$ is the bounded-check row, $O$ is the obstruction row, $H$ is $\hsame$ transport, $T$ is $\Cont$ replay, $G$ is $\Pkg$ provenance, and $N$ is the local $\NameCert_{\MetaCICParallelDiamondFrontierUp}$ row. Any confluence consumer must read this basis before it may cite the typed-window local join, and the basis does not discharge residual substitution, subject reduction, full closed Church--Rosser, evaluator equality, quotient equality of terms, or kernel consistency.
\end{theorem}
\begin{proof}
Start with the residual budget theorem \autoref{thm:metacic-parallel-diamond-frontier-budget}. It keeps the retained premise, critical-pair row, residual-frontier row, closed-checker row, closed-normal row, obstruction row, and structural rows visible before any exported budget is consumed.

Apply residual confluence and residual-join readback, \autoref{thm:metacic-parallel-diamond-frontier-residual-confluence} and \autoref{thm:metacic-parallel-diamond-frontier-residual-join-readback}. These route confluence-facing reads through the same residual spine and require bounded-conversion readback to keep the typed-window and L10 source rows visible.

Use the critical-overlap budget and candidate-mediated residual route, \autoref{thm:metacic-parallel-diamond-frontier-critical-overlap-budget} and \autoref{thm:metacic-parallel-diamond-frontier-candidate-mediated-residual-route}. The local peak may reach a candidate only through the displayed $K,J,R$ rows, with $P$ and $O$ retained.

Finally apply the typed-window local-join scope theorem \autoref{thm:metacic-parallel-diamond-frontier-local-join-scope}. It identifies the permitted local join as exactly the typed-window row surface, so the listed rows form the finite obligation basis and no stronger confluence endpoint is exported.
\end{proof}

\begin{theorem}[MetaCIC parallel-diamond frontier seed status]
\label{thm:metacic-parallel-diamond-frontier-seed-status}
Every conditional parallel-diamond handoff consumed through an accepted $\MetaCICParallelDiamondFrontierUp$ packet must first expose the residual-substitution premise row $P$, critical-pair row $K$, candidate-join row $J$, residual-frontier row $R$, checker row $S$, bounded-check row $B$, obstruction row $O$, componentwise $\hsame$ transport row $H$, displayed $\Cont$ replay row $T$, $\Pkg$ provenance row $G$, and local $\NameCert_{\MetaCICParallelDiamondFrontierUp}$ row $N$. The consumed handoff is therefore a seed-status frontier: it records which residual premises and finite checks remain visible before use, and it is not a parameter-transport equivalence, source-equivalence theorem, full Church--Rosser theorem, subject-reduction theorem, evaluator equality, quotient equality of terms, or kernel-consistency certificate.
\end{theorem}
\begin{proof}
Begin with the obligations theorem \autoref{thm:metacic-parallel-diamond-frontier-obligations}. It makes the residual-substitution premise, critical-pair, candidate-join, residual-frontier, checker, bounded-check, obstruction, transport, replay, provenance, and naming rows part of the finite carrier surface.

Apply the obligation-basis theorem \autoref{thm:metacic-parallel-diamond-frontier-obligation-basis}. It packages those rows as the typed-window basis that a confluence consumer must read before citing a local join.

The residual-frontier and obstruction rows stay visible throughout the handoff. Consequently the theorem records seed status for the remaining residual-discharge work rather than converting the packet into a parameter-transport equivalence or any stronger confluence endpoint.
\end{proof}

\begin{theorem}[MetaCIC parallel-diamond residual obstruction ledger]
\label{thm:metacic-parallel-diamond-frontier-residual-obstruction-ledger}
Every candidate-mediated local join read from an accepted $\MetaCICParallelDiamondFrontierUp$ packet must keep the residual-substitution premise row $P$, residual-frontier row $R$, and obstruction row $O$ visible alongside
$$
K,J,S,C,B,H,T,G,N
$$
before any confluence-facing packet consumes the read. The retained ledger is therefore
$$
P,K,J,R,S,C,B,O,H,T,G,N .
$$
It is not a premise-free parallel diamond, full Church--Rosser theorem, subject-reduction theorem, evaluator equality, quotient equality of terms, or kernel-consistency certificate.
\end{theorem}
\begin{proof}
Begin with the obligation basis \autoref{thm:metacic-parallel-diamond-frontier-obligation-basis}. It fixes the consumer-visible packet as $P,K,J,R,S,C,B,O,H,T,G,N$ and makes the obstruction row an explicit coordinate of the finite basis.

Apply the candidate-mediated residual route \autoref{thm:metacic-parallel-diamond-frontier-candidate-mediated-residual-route}. A local peak may reach the candidate row $J$ only through the displayed critical-pair and residual-frontier route, while $P$ remains the retained residual-substitution premise.

Finally use the local-join scope theorem \autoref{thm:metacic-parallel-diamond-frontier-local-join-scope}. It allows confluence-facing consumers to read only the typed-window surface with $O$ still visible, so the obstruction ledger cannot be erased before the handoff is consumed.

Thus every candidate-mediated join read carries $P,R,O$ together with $K,J,S,C,B,H,T,G,N$. Since neither the premise nor the obstruction is discharged, the statement exports only a residual obstruction ledger and not any refused confluence, typing, equality, quotient, or kernel endpoint.
\end{proof}

\begin{theorem}[MetaCIC residual pasting frontier]
\label{thm:metacic-parallel-diamond-frontier-residual-pasting}
Suppose an accepted $\MetaCICParallelDiamondFrontierUp$ packet has two typed-window residual reads whose critical-pair observations share the same $P,K,R,S,C,B,O,H,T,G,N$ rows and whose candidate row is $J$. Then the two reads paste only through the displayed residual packet
$$
J,R,S,C,B,O,H,T,G,N
$$
under the retained premise $P$. The pasted read is a finite $\NameCert$ route with $\hsame$ transport, $\Cont$ replay, and $\Pkg$ provenance; it is not a source-equivalence transport theorem, abstract carrier equivalence, premise-free parallel diamond, subject-reduction theorem, evaluator equality, quotient equality of terms, or kernel-consistency certificate.
\end{theorem}
\begin{proof}
Begin with \autoref{thm:metacic-parallel-diamond-frontier-obligation-basis}. It fixes the consumer-visible basis as $P,K,J,R,S,C,B,O,H,T,G,N$ and keeps the obstruction row visible.

Use candidate-mediated residual routing, \autoref{thm:metacic-parallel-diamond-frontier-candidate-mediated-residual-route}. The two local reductions may meet only at the displayed candidate row $J$ and residual-frontier row $R$, while the residual-substitution premise $P$ remains retained.

Apply residual confluence and residual-join readback. These theorems keep the closed checker, closed-normal fragment, bounded-check row, obstruction, transport, replay, provenance, and local naming attached to the same packet before any bounded-conversion consumer can read the handoff.

The pasted route is therefore the shared residual packet $J,R,S,C,B,O,H,T,G,N$ under $P$. Since the construction never introduces a source-equivalence carrier or removes $O$, the conclusion is a finite frontier pasting claim and not any refused confluence, typing, equality, quotient, or kernel endpoint.
\end{proof}

\begin{theorem}[MetaCIC substitution nonescape frontier]
\label{thm:metacic-parallel-diamond-frontier-substitution-nonescape}
In every accepted $\MetaCICParallelDiamondFrontierUp$ packet, substitution data remains confined to the residual-substitution boundary row $P$ and the residual-frontier row $R$ throughout the typed-window local join. No $\hsame$ transport row, $\Cont$ replay row, $\Pkg$ provenance row, or local $\NameCert_{\MetaCICParallelDiamondFrontierUp}$ row may reinterpret the handoff as a discharged residual-substitution theorem, premise-free parallel diamond, subject reduction, full closed Church--Rosser, evaluator equality, quotient equality of terms, or kernel consistency.
\end{theorem}
\begin{proof}
Start from the obligation basis \autoref{thm:metacic-parallel-diamond-frontier-obligation-basis}. It makes $P$ and $R$ explicit coordinates of the finite packet rather than implicit assumptions around the local join.

Apply the residual-substitution scope theorem and premise preservation, \autoref{thm:metacic-parallel-diamond-frontier-residual-substitution-scope} and \autoref{thm:metacic-parallel-diamond-frontier-premise-preservation}. They keep the substitution premise visible through the critical-pair, residual-frontier, closed-checker, and closed-normal rows.

Use residual pasting, \autoref{thm:metacic-parallel-diamond-frontier-residual-pasting}. Even when two typed-window residual reads are pasted through the candidate row, the pasted route still runs under $P$ and through $R$.

Finally read the structural rows $H,T,G,N$. They transport, replay, source, and name the displayed packet, but they do not contain a carrier slot for discharging substitution or erasing the obstruction row. Thus substitution cannot escape the displayed residual boundary into any stronger confluence endpoint.
\end{proof}

\begin{theorem}[MetaCIC parallel-diamond residual-substitution route]
\label{thm:metacic-parallel-diamond-residual-substitution-route}
Every candidate-mediated local join read from an accepted $\MetaCICParallelDiamondFrontierUp$ carrier $F=(P,K,J,R,S,C,B,O,H,T,G,N)$ must expose the residual-substitution premise row $P$, critical-pair row $K$, candidate-join row $J$, residual-frontier row $R$, closed-checker seal row $S$, closed-normal fragment row $C$, bounded-check row $B$, obstruction row $O$, componentwise $\hsame$ transport, displayed $\Cont$ replay, $\Pkg$ provenance, and local naming before any parallel-diamond frontier handoff is consumed. The route is conditional on the displayed premise and cannot be read as residual-substitution compatibility, full Church--Rosser, subject reduction, evaluator equality, quotient equality of terms, or kernel consistency.
\end{theorem}
\begin{proof}
Project the carrier of \autoref{def:metacic-parallel-diamond-frontier-carrier}. The rows $P,K,J,R,S,C,B,O,H,T,G,N$ are the full consumer-visible packet, with $P$ retaining the residual-substitution premise and $O$ retaining the refused endpoints.

Apply the candidate-mediated residual route of \autoref{thm:metacic-parallel-diamond-frontier-candidate-mediated-residual-route}. It allows the two local reductions named by $K$ to meet only through the displayed candidate row $J$ and residual-frontier row $R$, under the retained premise $P$.

Next use substitution nonescape, \autoref{thm:metacic-parallel-diamond-frontier-substitution-nonescape}. The closed checker row, closed-normal row, bounded-check row, transport, replay, provenance, and local naming cannot discharge substitution or remove the obstruction row.

Thus the route exposes exactly the residual-substitution premise, critical-pair, candidate-join, residual-frontier, closed-checker, closed-normal, bounded-check, obstruction, and structural rows stated above. Since neither $P$ nor $O$ disappears, the handoff remains a conditional frontier and not any refused confluence, typing, equality, quotient, or kernel endpoint.
\end{proof}

\begin{theorem}[MetaCIC parallel-diamond normalization-frontier handoff]
\label{thm:metacic-parallel-diamond-normalization-frontier-handoff}
Any normalization consumer that reads the $\MetaCICParallelDiamondFrontierUp$ packet may use only the conditional frontier exposed by \autoref{thm:metacic-parallel-diamond-residual-substitution-route}. It may inspect the retained residual-substitution premise, critical-pair row, candidate-join row, residual-frontier row, closed-checker row, closed-normal fragment row, bounded-check row, obstruction row, transport, replay, provenance, and local naming. It may not consume the packet as full Church--Rosser, subject reduction, evaluator equality, quotient equality of terms, kernel consistency, or a completed residual-substitution theorem.
\end{theorem}
\begin{proof}
Begin with the residual-substitution route theorem. It fixes the normalization-facing read to the displayed finite spine $P,K,J,R,S,C,B,O,H,T,G,N$ and keeps the obstruction row visible.

Use the typed-window local-join scope of \autoref{thm:metacic-parallel-diamond-frontier-local-join-scope}. A normalization consumer may read local joins only under that typed-window surface, so it cannot erase the residual-frontier or closed-checker conditions.

Apply residual pasting and substitution nonescape, \autoref{thm:metacic-parallel-diamond-frontier-residual-pasting} and \autoref{thm:metacic-parallel-diamond-frontier-substitution-nonescape}. These theorems keep pasted residual reads under $P$ and prevent structural rows from retyping the handoff as a discharged substitution result.

Consequently the normalization consumer receives only the conditional frontier. The listed stronger endpoints require separate certificates and are not carried by the accepted $\MetaCICParallelDiamondFrontierUp$ packet.
\end{proof}

\begin{theorem}[MetaCIC parallel-diamond candidate-mediated strong-normalization socket]
\label{thm:metacic-parallel-diamond-frontier-candidate-mediated-sn-socket}
For every accepted $\MetaCICParallelDiamondFrontierUp$ carrier $F=(P,K,J,R,S,C,B,O,H,T,G,N)$, suppose the closed checker row $S$ accepts both residual branches named by the local critical-pair row $K$, the candidate-normalization scope row $J$ retains the residual-substitution premise row $P$, and the local join ledger read through $R,C,B,H,T,G,N$ is available. Then the packet exposes a candidate-mediated strong-normalization socket for that local pair: its readback is the conditional route $P,K,J,R,S,C,B,O,H,T,G,N$ consumed by \autoref{thm:metacic-parallel-diamond-normalization-frontier-handoff}. The socket does not assert global strong normalization, full confluence, subject reduction, evaluator equality, quotient equality of terms, or kernel consistency.
\end{theorem}
\begin{proof}
Start from the critical-pair envelope of \autoref{thm:metacic-parallel-diamond-critical-pair-exhaustion}. It fixes the two residual branches inside the displayed local pair row $K$ and keeps them tied to the closed checker row $S$.

Use the residual-substitution dependency route of \autoref{thm:metacic-parallel-diamond-frontier-residual-substitution-scope}. The premise row $P$ remains an explicit coordinate of the candidate-normalization scope instead of becoming a discharged theorem.

Next apply the closed-normal join socket and local join ledger pullback, \autoref{thm:metacic-parallel-diamond-frontier-closed-normal-join-socket} and \autoref{thm:metacic-parallel-diamond-frontier-local-join-scope}. These results route the accepted branches through $J,R,C,B$ while preserving $H,T,G,N$ as transport, replay, provenance, and naming rows.

Finally hand the same finite readback to \autoref{thm:metacic-parallel-diamond-normalization-frontier-handoff}. Since $P$ and $O$ remain visible, the socket is only the candidate-mediated local normalization route promised in the statement, not any global normalization or confluence endpoint.
\end{proof}

\begin{theorem}[MetaCIC parallel-diamond frontier typed-window candidate SN route]
\label{thm:metacic-parallel-diamond-frontier-typed-window-candidate-sn-route}
Every strong-normalization-facing read of the candidate-mediated socket in an accepted $\MetaCICParallelDiamondFrontierUp$ carrier $F=(P,K,J,R,S,C,B,O,H,T,G,N)$ must pass through the typed-window handoff, residual-substitution scope, premise preservation, local-join scope, and the candidate-mediated strong-normalization socket before it reaches the normalization-frontier handoff. The resulting route is the finite readback
$$
P,K,J,R,S,C,B,O,H,T,G,N
$$
with the typed window still visible. It does not assert global strong normalization, full confluence, subject reduction, evaluator equality, quotient equality of terms, or kernel consistency.
\end{theorem}
\begin{proof}
Start with the typed-window handoff of \autoref{thm:metacic-parallel-diamond-frontier-typed-window-handoff}. It forces every candidate-normalization consumer to expose the same residual spine before any frontier handoff is read.

Next apply residual-substitution scope and premise preservation, \autoref{thm:metacic-parallel-diamond-frontier-residual-substitution-scope} and \autoref{thm:metacic-parallel-diamond-frontier-premise-preservation}. These keep $P$ attached to the critical-pair, residual-frontier, closed-checker, and closed-normal rows rather than turning it into a discharged substitution theorem.

Use the typed-window local-join scope of \autoref{thm:metacic-parallel-diamond-frontier-local-join-scope}. A candidate join can be consumed only through the displayed $K,J,R,S,C,B$ rows, with the obstruction row $O$ still retained.

Finally apply the candidate-mediated strong-normalization socket and normalization-frontier handoff, \autoref{thm:metacic-parallel-diamond-frontier-candidate-mediated-sn-socket} and \autoref{thm:metacic-parallel-diamond-normalization-frontier-handoff}. They pass the same finite readback forward as a conditional socket, so the route cannot be retyped as any of the refused normalization, confluence, typing, equality, quotient, or kernel endpoints.
\end{proof}

\begin{theorem}[MetaCIC parallel-diamond closed-substitution window scope]
\label{thm:metacic-parallel-diamond-frontier-closed-substitution-window-scope}
Every candidate-mediated diamond read of an accepted $\MetaCICParallelDiamondFrontierUp$ carrier $F=(P,K,J,R,S,C,B,O,H,T,G,N)$ that exposes the rows $P,K,J,R,S,C,B,O,H,T,G,N$ must keep the closed-substitution window visible before the closed checker seal and normalization-confluence handoff are consumed. The read therefore remains a finite frontier route through the displayed residual-substitution premise, typed candidate window, closed checker, local join ledger, obstruction row, transport, replay, provenance, and local naming rows. It cannot be reread as premise-free subject reduction or as closed Church--Rosser.
\end{theorem}
\begin{proof}
Project the carrier and begin with the premise-retention route of \autoref{thm:metacic-parallel-diamond-frontier-premise-preservation}. The residual-substitution premise row $P$ remains visible whenever the candidate rows $K,J,R,S,C,B,O$ are exposed.

Apply the typed-window candidate SN route of \autoref{thm:metacic-parallel-diamond-frontier-typed-window-candidate-sn-route}. It sends every strong-normalization-facing read through the typed window and candidate-mediated socket before the normalization-frontier handoff is available.

Finally use the critical-path retained-premise handoff of \autoref{thm:metacic-parallel-diamond-frontier-critical-path-retained-premise-handoff}. The structural rows $H,T,G,N$ transport, replay, source, and name the same retained-premise packet, so the closed checker seal and normalization-confluence handoff cannot erase the closed-substitution window. Hence the route is not a premise-free subject-reduction theorem and not closed Church--Rosser.
\end{proof}

\begin{theorem}[MetaCIC parallel-diamond finite candidate SN extension]
\label{thm:metacic-parallel-diamond-frontier-finite-candidate-sn-extension}
Let $F=(P,K,J,R,S,C,B,O,H,T,G,N)$ be an accepted $\MetaCICParallelDiamondFrontierUp$ carrier. If a finite local peak is routed through the displayed candidate row $J$, residual row $R$, typed window $S$, and boundary rows $C,B,O$, then the only strong-normalization-facing read is the candidate-mediated socket already named by \autoref{thm:metacic-parallel-diamond-frontier-candidate-mediated-sn-socket}. The route cannot be retyped as global strong normalization, full confluence, quotient equality of terms, or a kernel endpoint.
\end{theorem}
\begin{proof}
Apply the candidate-mediated residual route \autoref{thm:metacic-parallel-diamond-frontier-candidate-mediated-residual-route}. It sends the finite local peak through the displayed candidate row $J$ and residual row $R$ while retaining the residual-substitution premise row $P$.

Use the local-join scope \autoref{thm:metacic-parallel-diamond-frontier-local-join-scope}. The local join is visible only through the typed-window and boundary rows $S,C,B,O$, so no ambient confluence carrier is introduced.

Next apply the normalization-frontier handoff \autoref{thm:metacic-parallel-diamond-normalization-frontier-handoff}. It restricts normalization consumers to the conditional frontier packet and keeps the refused endpoints outside the carrier.

Finally use the candidate-mediated strong-normalization socket and typed-window candidate SN route, \autoref{thm:metacic-parallel-diamond-frontier-candidate-mediated-sn-socket} and \autoref{thm:metacic-parallel-diamond-frontier-typed-window-candidate-sn-route}. These results identify the only SN-facing read as the finite candidate-mediated socket through $P,K,J,R,S,C,B,O,H,T,G,N$, so it cannot be reread as global strong normalization, full confluence, quotient equality, or a kernel endpoint.
\end{proof}

\begin{theorem}[MetaCIC parallel-diamond frontier seed closure status]
\label{thm:metacic-parallel-diamond-frontier-seed-closure-status}
The seed closure surface of an accepted $\MetaCICParallelDiamondFrontierUp$ packet is exactly the finite conjunction of residual-substitution premise retention, critical-pair envelope exposure, closed-checker retention, closed-normal join socket, candidate-mediated strong-normalization socket, componentwise $\hsame$ transport, displayed $\Cont$ replay, $\Pkg$ provenance, and local $\NameCert$ rows. The seed surface does not discharge residual-substitution compatibility, premise-free parallel diamond, subject reduction, full closed Church--Rosser, evaluator equality, quotient equality of terms, or kernel consistency.
\end{theorem}
\begin{proof}
Start with residual-substitution scope and premise preservation, \autoref{thm:metacic-parallel-diamond-frontier-residual-substitution-scope} and \autoref{thm:metacic-parallel-diamond-frontier-premise-preservation}. They keep the row $P$ visible through every local-diamond read.

Use the critical-pair envelope and closed-checker retention, \autoref{thm:metacic-parallel-diamond-critical-pair-exhaustion} and \autoref{thm:metacic-parallel-diamond-closed-checker-retention}. These identify the rows $K,S,C,B,O$ as part of the finite frontier instead of consequences of a discharged confluence theorem.

Apply the local-join and normalization handoff theorems, \autoref{thm:metacic-parallel-diamond-frontier-local-join-scope} and \autoref{thm:metacic-parallel-diamond-normalization-frontier-handoff}. Together with the candidate-mediated strong-normalization socket \autoref{thm:metacic-parallel-diamond-frontier-candidate-mediated-sn-socket}, they show that the only normalization-facing read is the conditional route through $P,K,J,R,S,C,B,O$.

Finally read the structural rows $H,T,G,N$. They transport, replay, source, and name the same finite packet, so the seed closure surface is exactly the listed conjunction and none of the refused endpoints is exported.
\end{proof}

\begin{theorem}[MetaCIC parallel-diamond residual-join induction]
\label{thm:metacic-parallel-diamond-frontier-residual-join-induction}
Let $F=(P,K,J,R,S,C,B,O,H,T,G,N)$ be an accepted $\MetaCICParallelDiamondFrontierUp$ carrier. Finite iteration over a displayed list of residual pairs preserves the local-diamond route through the critical-pair row $K$, residual-substitution premise row $P$, candidate-normalization row $J$, closed-checker row $S$, bounded-check row $B$, obstruction row $O$, transport row $H$, replay row $T$, provenance row $G$, and local naming row $N$. The iterated route remains candidate-mediated and exports no premise-free parallel diamond, subject reduction, full closed Church--Rosser theorem, evaluator equality, quotient equality of terms, or kernel consistency endpoint.
\end{theorem}
\begin{proof}
Induct on the finite residual-pair list. The base case is \autoref{thm:metacic-parallel-diamond-frontier-seed-closure-status}: seed readiness already exposes residual-substitution premise retention, critical-pair rows, candidate joins, closed checker rows, bounded checks, obstruction retention, and the structural rows.

For the step, append one displayed residual pair. The residual-substitution dependency route of \autoref{thm:metacic-parallel-diamond-frontier-residual-substitution-scope} keeps $P$ visible, while candidate-normalization scope and finite residual witness consumption keep the new pair inside $K,J,R,S,C,B,O$ before any normalization-facing handoff is read.

Finally apply the critical-path consumer route of \autoref{thm:metacic-parallel-diamond-frontier-critical-path-retained-premise-handoff}. It passes the extended finite readback through $H,T,G,N$ without changing the refused endpoints, so the induction preserves the local-diamond route claimed above.
\end{proof}

The Lean follow-up supplies $\mathsf{Nontrivial}$ and $\mathsf{FieldFaithful}$ TasteGate instances in \texttt{lean4/BEDC/Derived/MetaCICParallelDiamondFrontierUp/TasteGate.lean} by exhibiting distinct accepted and refused frontier packets and by matching on the rows $P,K,J,R,S,C,B,O,H,T,G,N$.

\closureat{MetaCICParallelDiamondFrontierUp}{scopedStr}

\begin{closurestatus}{\MetaCICParallelDiamondFrontierUp}
  \constructivestory{A $\MetaCICParallelDiamondFrontierUp$ packet is generated from residual-substitution premise rows, critical-pair rows, candidate-join rows, residual-frontier rows, closed checker seal rows, closed-normal fragment rows, bounded-check rows, obstruction rows, componentwise hsame transport, displayed Cont replay, Pkg provenance, and a local NameCert row.}
  \theoryclosure{\scopedClosure}
  \scopeclosed{\autoref{def:metacic-parallel-diamond-frontier-carrier}, \autoref{thm:metacic-parallel-diamond-frontier-residual-substitution-scope}, \autoref{thm:metacic-parallel-diamond-critical-pair-exhaustion}, \autoref{thm:metacic-parallel-diamond-frontier-obstruction-triad}, \autoref{thm:metacic-parallel-diamond-frontier-residual-peak-scope}, \autoref{thm:metacic-parallel-diamond-closed-checker-retention}, \autoref{thm:metacic-parallel-diamond-frontier-local-join-scope}, \autoref{thm:metacic-parallel-diamond-normalization-frontier-handoff}, \autoref{thm:metacic-parallel-diamond-frontier-candidate-mediated-sn-socket}, and \autoref{thm:metacic-parallel-diamond-frontier-residual-join-induction} delimit the residual-premise, critical-pair, candidate-join, residual-frontier, closed-checker, closed-normal, bounded-check, obstruction, residual-join induction, and candidate-mediated strong-normalization rows for MetaCIC confluence consumers.}
  \formalstatus{\unformalizedV}
  \bridgestatus{none}
  \notclaimed{This chapter does not close residual-substitution compatibility, full parallel diamond, subject reduction, arbitrary closed Church--Rosser, evaluator equality, quotient equality of terms, kernel consistency, Lean verification, or bridge checking.}
  \upgradepath{Obligation closure requires checked carrier admission, residual-substitution premise retention, critical-pair exposure, candidate-join exactness, residual-frontier coverage, closed-checker handoff, closed-normal fragment routing, bounded-check exactness, obstruction retention, TasteGate nontriviality, FieldFaithful coverage, transport, replay, provenance, and local naming for BEDC.Derived.MetaCICParallelDiamondFrontierUp.}
\end{closurestatus}
